\section{Quantitative accounting for a future hierarchy engine}
\label{sec:hierarchy-accounting}

The near-linear terminal test removes one repeated search from the
hierarchy route.  To obtain a general quasi-polynomial recognizer, one
must also bound the construction and revision of the hierarchy, the
size of its compressed states, and the cost of transporting a terminal
witness.  This section specifies a sufficient global accounting theorem.
It does not assert that the geometric hypotheses have been established
for the delivered implementation.

\subsection{The precise status of the source bounds}

The February 2021 slides cited here, the file
\code{quasipolynomial-talk.pdf}, announce
$n^{O(\log n)}$ \cite{lackenbytalk}.  The November 2024 Tsinghua abstract
announces
\[
 2^{O((\log n)^3)}=n^{O((\log n)^2)}
 \qquad\cite{lackenby2024}.
\]
These are different stated bounds.  Both are quasi-polynomial, and the
sources reviewed here do not establish an explanation for the difference.
The July 2026 preprint is a separate source: it gives a hierarchy
algorithm and efficient certificate operations, while identifying further
speed-up as future work \cite{lackenby2026}.

For its iteration analysis, let $L$ bound the length of every hierarchy
that occurs, and let $g$ bound every nonnegative complexity digit that
is stored, including after a suffix is rebuilt.  Proposition~9.1 gives
\begin{equation}
 \#\{\hbox{executions of any fixed algorithmic step}\}
 \leq L(g+1)^L.
 \label{eq:lackenby-iteration-bound}
\end{equation}
Proposition~10.2 gives the refined length bound
$L\leq4c_q(\mathcal H)$, in terms of the initial handle structure's
q-complexity.  For a linear-size initial structure of uniform type this
is $L=O(n)$ \cite[Props.~9.1 and 10.2]{lackenby2026}.

Equation~\eqref{eq:lackenby-iteration-bound} is not yet a bit-time estimate.
Its digits concern pattern complexity, including intersections with the
boundary pattern, rather than genus alone.  A bound valid only for the
first chosen surfaces is insufficient when later replacements introduce
new surfaces.  Moreover, each counted step may act on a compressed object
whose bit length must itself be controlled.

Even if the additional estimate $g\leq(n+2)^a$ were supplied, combining it
only with $L=O(n)$ would give the coarse iteration bound
$\exp(O(n\log(n+2)))$.  This is not a quasi-polynomial upper bound.
It is also not a lower bound on the algorithm: a sharper analysis may
exploit correlations between successive choices or long stretches of
forced operations.  The distinction motivates the following elementary
accounting result.

\subsection{A rooted-tree bound that includes forced operations}

All logarithms in this subsection are natural logarithms.  A rooted tree
represents \emph{actual execution occurrences}: distinct visits are
distinct nodes unless a separate memoization argument proves that they
can be charged to one computation.  At a node $v$, write $b(v)$ for the
number of its children.  A node with one child describes a forced
operation and contributes $\log b(v)=0$.  Leaves contribute no logarithm.

\begin{theorem}[Depth and cumulative branching]
\label{thm:hierarchy-tree-accounting}
Let $T$ be a rooted, finitely branching tree of depth at most $D$, where
depth counts edges.  Suppose that every root-to-leaf path satisfies
\begin{equation}
 \sum_{\substack{v\ \mathrm{on\ the\ path}\\b(v)>0}}\log b(v)\leq B.
 \label{eq:branch-budget}
\end{equation}
Then
\begin{equation}
 |V(T)|\leq(D+1)e^B.
 \label{eq:tree-node-bound}
\end{equation}
If all work associated with one node has bit cost at most $C$, including
its assigned state construction, tests, certificate data, and transition
work, the total assigned cost is at most $C(D+1)e^B$.
\end{theorem}

\begin{proof}
Finite branching and bounded depth imply that $T$ is finite.  Assign the
root weight $1$, and assign every child of an internal node $v$ the
weight of $v$ divided by $b(v)$.  Thus at every fixed depth the sum of
the node weights is at most $1$: passing to the next depth preserves the
weight at internal nodes and drops the weight at leaves.

For a node $u$, its weight is
\[
 p(u)=
 \prod_{\substack{v\ \mathrm{a\ proper}\\\mathrm{ancestor\ of}\ u}}
 \frac1{b(v)}.
\]
Extend its path to a leaf.  Since every summand in
\eqref{eq:branch-budget} is nonnegative, the ancestor sum for $u$ is at
most $B$.  Hence $p(u)\geq e^{-B}$.  If there are $N_j$ nodes at depth
$j$, then $N_je^{-B}\leq1$, so $N_j\leq e^B$.  Summing over
$j=0,\ldots,D$ proves~\eqref{eq:tree-node-bound}.  Summing the stated
per-node costs proves the final assertion.
\end{proof}

The weights are a counting device.  No randomized algorithm or
probabilistic running-time assumption is involved.  The same proof works
with computable bounds $\widehat b(v)\geq b(v)$ in place of the actual
positive child counts: distribute weight $p(v)/\widehat b(v)$ to each
child, so that total weight can only decrease.

Unary chains explain why the depth term cannot be omitted.  A chain of
$D+1$ nodes has $B=0$ and attains the bound exactly.  Conversely, a
small bound on the length of each completed geometric hierarchy is not
automatically a bound on $D$: rebuilding its suffix many times can
create many execution occurrences.  A tree interpretation must account
for those occurrences explicitly.

\begin{corollary}[Explicit quasi-polynomial budgets]
\label{cor:hierarchy-qp-budget}
Suppose a complete algorithm on $n$-crossing inputs admits the accounting
of Theorem~\ref{thm:hierarchy-tree-accounting}, uniformly over all its
inputs and executions.  If
\[
 B(n)+\log(D(n)+1)+\log(C(n)+1)
     =O\bigl((\log(n+2))^r\bigr)
\]
for a fixed $r\geq1$, then its accounted bit time is
$\exp(O((\log(n+2))^r))$.
In particular, polynomial $C,D$ and
$B=O(\log^2(n+2))$ imply $n^{O(\log(n+2))}$ time; replacing that
branching budget by $O(\log^3(n+2))$ gives
$n^{O((\log(n+2))^2)}$.
\end{corollary}

\begin{proof}
Take logarithms of the time bound in
Theorem~\ref{thm:hierarchy-tree-accounting} and substitute the hypotheses.
\end{proof}

For example, polynomially many forced operations along a path are
compatible with the stronger bound when at most $O(\log n)$ nodes
branch and each has at most $n^{O(1)}$ children.  Their cumulative
branching is $O(\log^2 n)$.  If polynomial branching occurs at
$O(\log^2 n)$ levels instead, the same arithmetic gives
$O(\log^3 n)$.  These are sufficient accounting regimes, not an
explanation of the historical announcements.

This formulation permits the available choices to depend on the entire
earlier history.  It can therefore be more informative than multiplying
independent worst-case caps at every hierarchy level.  Nevertheless,
the bound must hold on all explored paths, including failed choices and
restarts.  A bound on one successful certificate, or on the number of
distinct unlabeled states, does not establish the hypotheses.

\subsection{The geometric obligations behind the symbols}

There are three distinct representation issues.  First, a polynomial
number of local handle types does not bound the number of handles.
Second, large normal coordinates require binary arithmetic and
multiplicity-aware output, as
Proposition~\ref{prop:parallel-disc-output} demonstrates.  Third, a
compact object can still require a costly search to construct.
The per-node bound $C$ must include all three issues rather than treating
surface production or witness transport as an oracle.

Theorem~14.4 of the 2026 paper supplies compressed cutting data in time
polynomial in the handle count and logarithm of the extended surface
weight.  Proposition~14.2's subsequent simplification has additional
irreducibility and essential-pattern hypotheses
\cite[Prop.~14.2 and Thm.~14.4]{lackenby2026}.
The latter hypothesis cannot be silently supplied by an undecided knot
exterior with empty initial pattern.  A verifier may establish it from
an appropriate certificate; a general producer must also handle the
opposite case.

A concrete realization of Corollary~\ref{cor:hierarchy-qp-budget} should
therefore attach to every state an encoding-size bound, an explicit
description of its admissible continuations, and a record of each
operation used to transfer boundary curves.  If a simplification is
conditional on essentiality, its contract must state how that
precondition is certified or how a compression witness is handled.
All outer changes of Heegaard data and all reconstructions of discarded
suffixes must receive a charge in the same execution accounting.

Efficient verification of a supplied hierarchy and efficient generation
of one are separate requirements.  The present package advances the
terminal test and the algebraic route, and gives exact interfaces against
which a hierarchy implementation can be audited.  It does not supply
uniform geometric bounds for $B,D,C$.  Proving those bounds, together
with an executable producer satisfying their hypotheses, remains the
step needed for a general hierarchy-based complexity theorem.
